\documentclass[a4paper]{article}
% Shared preamble for the labs (compile from Labs/ with XeLaTeX or tectonic).
% Same fonts and header as the Series sheets, plus colored boxes and tables.

\usepackage{enumitem}
\usepackage[margin=2.2cm]{geometry}
\usepackage{amsmath,amsfonts,amssymb}
\usepackage[charter]{mathdesign}
\usepackage{array}
\usepackage{colortbl}
\usepackage{fancyhdr}
\usepackage[skins,breakable]{tcolorbox}
\usepackage{polyglossia}

%------------------------------------------------
% Languages and fonts
%------------------------------------------------
\setdefaultlanguage[calendar=gregorian, numerals=maghrib]{arabic}
\setotherlanguage{english}

\IfFontExistsTF{Amiri}{%
    \newfontfamily\arabicfont[Script=Arabic, Scale=1.15]{Amiri}
}{%
    \newfontfamily\arabicfont[
        Script=Arabic, Scale=1.15,
        Path=../fonts/, Extension=.ttf,
        UprightFont=*-Regular, BoldFont=*-Bold,
        ItalicFont=*-Italic, BoldItalicFont=*-BoldItalic
    ]{Amiri}
}

\IfFontExistsTF{AlBattar}{%
    \newfontfamily\arabicfontsf[Script=Arabic, Scale=1.2]{AlBattar}
    \newfontfamily\khatt[Script=Arabic, Scale=1.2]{AlBattar}
}{%
    \newfontfamily\arabicfontsf[
        Script=Arabic, Scale=1.2, Path=../fonts/,
        BoldFont=ae_AlBattar.ttf
    ]{ae_AlBattar.ttf}
    \newfontfamily\khatt[
        Script=Arabic, Scale=1.2, Path=../fonts/,
        BoldFont=ae_AlBattar.ttf
    ]{ae_AlBattar.ttf}
}

%------------------------------------------------
% Colors (same palette as the slides)
%------------------------------------------------
\definecolor{brandTeal}{HTML}{00B4D8}
\definecolor{brandOrange}{HTML}{FF7043}
\definecolor{brandDark}{HTML}{263238}
\definecolor{brandMuted}{HTML}{6C757D}
\colorlet{tealDeep}{brandTeal!70!black}

%------------------------------------------------
% Page setup
%------------------------------------------------
\renewcommand{\footrulewidth}{1pt}
\renewcommand{\headrulewidth}{1pt}
\rhead{\khatt المدرسة العليا للأساتذة ببوزريعة}
\lhead{\khatt ملحقة بومرداس}
\cfoot{\thepage}
\lfoot{\khatt الأستاذ: د. بن عبد الرزاق يوسف}
\parindent 0pt
\setlength{\headsep}{15pt}
\setlength{\headheight}{16pt}
\setlength{\parskip}{4pt}
\pagestyle{fancy}

\setlist[enumerate,1]{label=\arabic*), itemsep=4pt, topsep=3pt}
\setlist[enumerate,2]{label=\alph*), itemsep=2pt, topsep=2pt}
\setlist[itemize]{itemsep=2pt, topsep=2pt}

%------------------------------------------------
% Tables
%------------------------------------------------
\renewcommand{\arraystretch}{1.4}
\setlength{\arrayrulewidth}{0.6pt}
\arrayrulecolor{brandMuted!80}
% Header cell
\newcommand{\ras}[1]{\cellcolor{brandTeal!20}\textbf{#1}}
% Highlighted result cell
\newcommand{\natija}[1]{\cellcolor{brandOrange!15}\textbf{#1}}
% Empty cell to be filled by the student
\newcommand{\fragh}{\rule{0pt}{0.55cm}}
\newcolumntype{C}[1]{>{\centering\arraybackslash}m{#1}}
\newcolumntype{R}[1]{>{\raggedleft\arraybackslash}m{#1}}

%------------------------------------------------
% Boxes
%------------------------------------------------
\tcbset{labbox/.style={enhanced, breakable, arc=3pt, boxrule=1.2pt,
    coltitle=white, fonttitle=\khatt\bfseries,
    top=1.5mm, bottom=1.5mm, left=3mm, right=3mm,
    before skip=6pt, after skip=6pt}}
% Worked example
\newtcolorbox{mithal}[1]{labbox, colback=brandTeal!6, colframe=brandTeal,
    title={مثال محلول: #1}}
% Reminder / method
\newtcolorbox{tadhkir}[1][تذكير]{labbox, colback=brandTeal!3,
    colframe=tealDeep, title={#1}}
% Warning
\newtcolorbox{tanbih}[1][انتبه]{labbox, colback=brandOrange!8,
    colframe=brandOrange, title={#1}}
% Exercise
\newtcolorbox{tamrin}[1]{labbox, colback=white, colframe=brandDark!80,
    title={#1}}

% Section title
\newcommand{\qism}[1]{\section*{\khatt\underline{#1}}}

% Title block: \unwan{lab number}{title}{subtitle}
\newcommand{\unwan}[3]{%
    \begin{center}
        {\Large\khatt\textbf{المخبر التطبيقي رقم #1}}\\[0.15cm]
        {\Large\khatt\textbf{#2}}\\[0.1cm]
        {\khatt\textbf{#3}}
    \end{center}
    \rule{\textwidth}{1pt}
    \begin{center}
        \small المدة: ساعة ونصف \quad$\cdot$\quad الأدوات: آلة حاسبة، أداة ذكاء اصطناعي، مصدر موثوق
    \end{center}
    \rule{\textwidth}{1pt}}

% Percentage reminder used in every lab
\newcommand{\nisba}{%
    \begin{tadhkir}[تذكير: حساب النسبة المئوية]
        \begin{center}
            \textbf{النسبة المئوية = (الجزء $\div$ الكل) $\times$ 100}
        \end{center}
        مثال: 15 إجابة صحيحة من 20 سؤالًا:
        $\dfrac{15}{20}\times 100 = 75\%$.
    \end{tadhkir}}


\begin{document}

\unwan{2}{كيف يتعلم الذكاء الاصطناعي؟}{البيانات ــ التدريب ــ الاختبار ــ الانحياز بالأرقام}

\qism{أهداف المخبر}
\begin{itemize}
    \item تقسيم مجموعة بيانات إلى بيانات تدريب وبيانات اختبار.
    \item حساب دقة نموذج على بيانات الاختبار.
    \item اكتشاف عدم توازن البيانات وأثره على النتائج.
    \item التمييز بين نموذج «فهم» ونموذج «حفظ».
    \item اكتشاف الانحياز الخوارزمي في جدول من الأرقام.
\end{itemize}

\begin{tanbih}[ملاحظة حول البيانات]
    الأرقام الواردة في هذا المخبر أمثلة تعليمية مبسطة، صُمّمت لتوضيح مفاهيم الدرس.
\end{tanbih}

\nisba

\begin{tadhkir}[تذكير: دقة النموذج]
    \begin{center}
        \textbf{الدقة = (عدد الإجابات الصحيحة $\div$ عدد أمثلة الاختبار) $\times$ 100}
    \end{center}
\end{tadhkir}

%================================================
\qism{الجزء الأول: أمثلة محلولة}
%================================================

\begin{mithal}{تقسيم البيانات وحساب الدقة}
جمع فريق \textbf{1000 جملة} عربية، وكتب أمام كل جملة نوعها: \textbf{اسمية} أو \textbf{فعلية}.
قسّم الفريق البيانات بنسبة \textbf{80\% للتدريب} و\textbf{20\% للاختبار}.

\begin{center}
\begin{tabular}{|R{4.5cm}|C{3cm}|C{3cm}|}
    \hline
    \ras{الجزء} & \ras{النسبة} & \ras{عدد الجمل} \\ \hline
    بيانات التدريب   & 80\%  & $1000 \times 0.8 = 800$ \\ \hline
    بيانات الاختبار  & 20\%  & $1000 \times 0.2 = 200$ \\ \hline
    \textbf{المجموع} & 100\% & \textbf{1000} \\ \hline
\end{tabular}
\end{center}

بعد التدريب، صنّف النموذج 200 جملة اختبار، فأصاب في \textbf{170} منها:
\begin{center}
    \textbf{الدقة} $= \dfrac{170}{200} \times 100 = 85\%$
\end{center}
\textbf{الاستنتاج:} النموذج يصيب في 85 جملة من كل 100،
ويخطئ في $200 - 170 = 30$ جملة.
\end{mithal}

\begin{mithal}{البيانات غير المتوازنة تخفي الضعف}
درّب طالب نموذجًا للتعرف على الأجهزة بهذه الصور، ثم اختبره بصور جديدة:

\begin{center}
\begin{tabular}{|R{2.4cm}|C{2.1cm}|C{2.4cm}|C{2.1cm}|C{1.7cm}|C{1.8cm}|}
    \hline
    \ras{الفئة} & \ras{صور التدريب} & \ras{نسبتها من التدريب} & \ras{صور الاختبار} & \ras{صحيحة} & \ras{الدقة} \\ \hline
    هواتف       & 500  & $45.5\%$ & 100 & 92  & 92\% \\ \hline
    حواسيب      & 500  & $45.5\%$ & 100 & 90  & 90\% \\ \hline
    ألواح رقمية & 100  & \natija{$9.1\%$} & 20  & 8   & \natija{40\%} \\ \hline
    \textbf{المجموع} & \textbf{1100} & 100\% & \textbf{220} & \textbf{190} & $86.4\%$ \\ \hline
\end{tabular}
\end{center}

\textbf{الحساب:}
\begin{itemize}
    \item نسبة الألواح من التدريب: $\dfrac{100}{1100}\times 100 \approx 9.1\%$.
    \item الدقة الإجمالية: $\dfrac{190}{220}\times 100 \approx 86.4\%$، وهي نتيجة تبدو جيدة.
    \item لكن دقة الألواح: $\dfrac{8}{20}\times 100 = 40\%$ فقط.
\end{itemize}
\textbf{الاستنتاج:} الدقة الإجمالية \textbf{أخفت} ضعف النموذج في الفئة القليلة.
لذلك نحسب الدقة لكل فئة على حدة، ونوازن بين الفئات في بيانات التدريب.
\end{mithal}

\begin{mithal}{نموذج فهم أم نموذج حفظ؟}
اختبر فريق نموذجين على بيانات التدريب نفسها، ثم على بيانات اختبار جديدة:

\begin{center}
\begin{tabular}{|R{3cm}|C{3cm}|C{3cm}|C{3cm}|}
    \hline
    \ras{النموذج} & \ras{دقة التدريب} & \ras{دقة الاختبار} & \ras{الفرق} \\ \hline
    أ & 99\% & 62\% & \natija{37 نقطة} \\ \hline
    ب & 88\% & 85\% & 3 نقاط \\ \hline
\end{tabular}
\end{center}

\textbf{الاستنتاج:} النموذج \textbf{أ} يشبه التلميذ الذي حفظ أمثلة القسم:
ممتاز في الأمثلة التي رآها، وضعيف في الأمثلة الجديدة.
النموذج \textbf{ب} فهم القاعدة، لأن نتيجته في الاختبار قريبة من نتيجته في التدريب.
نختار النموذج \textbf{ب}.
\end{mithal}

%================================================
\qism{الجزء الثاني: تمارين تطبيقية}
%================================================

\begin{tamrin}{التمرين 1: قسّم البيانات}
أكمل الجدول بحساب عدد أمثلة التدريب وعدد أمثلة الاختبار:

\begin{center}
\begin{tabular}{|C{3.2cm}|C{3.2cm}|C{3.2cm}|C{3.2cm}|}
    \hline
    \ras{عدد الأمثلة} & \ras{النسبة (تدريب / اختبار)} & \ras{التدريب} & \ras{الاختبار} \\ \hline
    2500 & 80 / 20 & 2000 & 500 \\ \hline
    1200 & 75 / 25 & \fragh & \\ \hline
    3000 & 70 / 30 & \fragh & \\ \hline
    450  & 90 / 10 & \fragh & \\ \hline
    640  & 80 / 20 & \fragh & \\ \hline
\end{tabular}
\end{center}
السطر الأول محلول مثالًا. لماذا يكون جزء التدريب دائمًا أكبر من جزء الاختبار؟
\end{tamrin}

\begin{tamrin}{التمرين 2: كن أنت النموذج}
هذه \textbf{بيانات التدريب}: ست جمل مصنّفة.

\begin{center}
\begin{tabular}{|C{0.6cm}|R{6cm}|C{3cm}|}
    \hline
    \ras{ر} & \ras{الجملة} & \ras{النوع} \\ \hline
    1 & كتب التلميذُ الدرسَ.         & فعلية \\ \hline
    2 & العلمُ نورٌ.                 & اسمية \\ \hline
    3 & ذهب الأستاذُ إلى القسم.      & فعلية \\ \hline
    4 & الكتابُ مفيدٌ.               & اسمية \\ \hline
    5 & يقرأ الطالبُ القصيدةَ.       & فعلية \\ \hline
    6 & السماءُ صافيةٌ.              & اسمية \\ \hline
\end{tabular}
\end{center}

\begin{enumerate}
    \item ما «النمط» الذي تكتشفه في هذه البيانات؟ اكتبه في شكل قاعدة واحدة.
    \item طبّق قاعدتك وحدها (كما تفعل الآلة) على \textbf{بيانات الاختبار}:

    \begin{center}
    \begin{tabular}{|C{0.6cm}|R{4.4cm}|C{2.6cm}|C{2.6cm}|}
        \hline
        \ras{ر} & \ras{الجملة} & \ras{توقع قاعدتك} & \ras{النوع الصحيح} \\ \hline
        أ & نجح المجتهدُ.            & \fragh & \\ \hline
        ب & الحديقةُ جميلةٌ.         & \fragh & \\ \hline
        ج & يكتب الشاعرُ قصيدةً.     & \fragh & \\ \hline
        د & إنّ العلمَ نافعٌ.         & \fragh & \\ \hline
        هـ & قد حضر الضيفُ.          & \fragh & \\ \hline
    \end{tabular}
    \end{center}
    \item احسب دقة قاعدتك على بيانات الاختبار.
    \item في أي جملة أخطأت القاعدة؟ ما الذي كان ناقصًا في بيانات التدريب؟
    \item اقترح جملتين تضيفهما إلى بيانات التدريب لتحسين «النموذج».
\end{enumerate}
\end{tamrin}

\begin{tamrin}{التمرين 3: نموذج تصحيح الأخطاء}
درّبت شركة نموذجًا يكشف الأخطاء اللغوية على هذه النصوص، ثم اختبرته بـ 100 نص من كل نوع.
تعني «الإجابة الصحيحة» أن النموذج لم يعدّ الكلمات الصحيحة أخطاءً.

\begin{center}
\begin{tabular}{|R{3.6cm}|C{2.4cm}|C{2.2cm}|C{2.4cm}|C{2.2cm}|}
    \hline
    \ras{نوع النص} & \ras{نصوص التدريب} & \ras{نسبتها} & \ras{صحيحة من 100} & \ras{الدقة} \\ \hline
    مقالات صحفية حديثة & 9000 & \fragh & 95 & \fragh \\ \hline
    روايات معاصرة       & 700  & \fragh & 88 & \fragh \\ \hline
    كتب مدرسية          & 250  & \fragh & 84 & \fragh \\ \hline
    شعر جاهلي           & 50   & \fragh & 41 & \fragh \\ \hline
    \textbf{المجموع}    & \fragh & 100\% & \fragh & \fragh \\ \hline
\end{tabular}
\end{center}

\begin{enumerate}
    \item أكمل الجدول.
    \item ما العلاقة بين عدد نصوص التدريب ودقة النموذج؟
    \item لماذا يخطئ النموذج كثيرًا في الشعر الجاهلي؟
    \item نريد على الأقل \textbf{1000 نص} من كل نوع. كم نصًا يجب أن نضيف لكل نوع؟
\end{enumerate}
\end{tamrin}

\begin{tamrin}{التمرين 4: هل النظام عادل؟}
درّبت شركة نظامًا يحوّل الكلام المنطوق إلى نص مكتوب، على تسجيلات صوتية من عدة لهجات جزائرية:

\begin{center}
\begin{tabular}{|R{3.4cm}|C{2.6cm}|C{3cm}|C{2.6cm}|}
    \hline
    \ras{اللهجة} & \ras{ساعات التسجيل} & \ras{نسبتها من التدريب} & \ras{دقة التعرف} \\ \hline
    العاصمية         & 400  & \fragh & 93\% \\ \hline
    الوهرانية        & 300  & \fragh & 90\% \\ \hline
    القسنطينية       & 250  & \fragh & 88\% \\ \hline
    لهجات الجنوب     & 50   & \fragh & 61\% \\ \hline
    \textbf{المجموع} & \fragh & 100\% & -- \\ \hline
\end{tabular}
\end{center}

\begin{enumerate}
    \item أكمل الجدول.
    \item ما الفرق بين أعلى دقة وأدنى دقة؟
    \item هل يمكن الحديث هنا عن \textbf{انحياز خوارزمي}؟ علّل بالأرقام.
    \item من هم المستخدمون الذين سيتضررون؟ اقترح حلين.
\end{enumerate}
\end{tamrin}

\begin{tamrin}{التمرين 5: اختر النموذج المناسب}
\begin{center}
\begin{tabular}{|C{2cm}|C{2.6cm}|C{2.6cm}|C{2.2cm}|R{3.4cm}|}
    \hline
    \ras{النموذج} & \ras{دقة التدريب} & \ras{دقة الاختبار} & \ras{الفرق} & \ras{الحكم} \\ \hline
    أ & 98\% & 64\% & \fragh & \\ \hline
    ب & 86\% & 83\% & \fragh & \\ \hline
    ج & 70\% & 68\% & \fragh & \\ \hline
\end{tabular}
\end{center}
\begin{enumerate}
    \item أكمل عمود الفرق.
    \item أيّ نموذج «حفظ» ولم يفهم؟ وأيّ نموذج لم يتعلم جيدًا أصلًا؟
    \item أيّ نموذج تختار؟ علّل.
\end{enumerate}
\end{tamrin}

%================================================
\newpage
\qism{الحلول المختصرة}
%================================================

\textbf{التمرين 1:}
1200: 900 / 300 \quad 3000: 2100 / 900 \quad 450: 405 / 45 \quad 640: 512 / 128.\\
جزء التدريب أكبر لأن النموذج يحتاج أمثلة كثيرة ليكتشف الأنماط؛ ويكفي جزء أصغر لتقييمه.

\textbf{التمرين 2:}
القاعدة: «إذا بدأت الجملة بفعل فهي فعلية، وإلا فهي اسمية».\\
التوقعات: أ) فعلية \quad ب) اسمية \quad ج) فعلية \quad د) اسمية \quad هـ) اسمية.\\
الصحيح: أ) فعلية \quad ب) اسمية \quad ج) فعلية \quad د) اسمية \quad هـ) \textbf{فعلية}.\\
الدقة: $\dfrac{4}{5}\times 100 = 80\%$.
أخطأت القاعدة في (هـ) لأن بيانات التدريب لم تحتوِ على جمل تبدأ بحرف مثل «قد» أو «إنّ».
أمثلة للإضافة: «قد نجح الطالبُ» (فعلية)، «لعلّ الجوَّ معتدلٌ» (اسمية).

\textbf{التمرين 3:}
المجموع: 10000 نص. النسب: $90\%$، $7\%$، $2.5\%$، $0.5\%$.
الدقة: $95\%$، $88\%$، $84\%$، $41\%$.
المجموع الصحيح: $95 + 88 + 84 + 41 = 308$ من 400، أي $77\%$.\\
كلما قلّت نصوص التدريب ضعفت الدقة. الشعر الجاهلي يمثل $0.5\%$ فقط من التدريب،
فلم يتعلم النموذج مفرداته القديمة.
الإضافة: روايات 300، كتب مدرسية 750، شعر جاهلي 950 (الصحفية لا تحتاج إضافة).

\textbf{التمرين 4:}
المجموع 1000 ساعة. النسب: $40\%$، $30\%$، $25\%$، $5\%$.
الفرق: $93 - 61 = 32$ نقطة.
نعم، انحياز: الفئة الممثلة بـ $5\%$ فقط من البيانات حصلت على أضعف دقة بفارق كبير.
المتضررون: سكان الجنوب. الحلول: جمع تسجيلات أكثر من لهجات الجنوب،
واختبار النظام على كل لهجة على حدة قبل نشره.

\textbf{التمرين 5:}
الفرق: أ $34$ نقطة، ب $3$ نقاط، ج $2$ نقطتان.
أ «حفظ» (فرق كبير). ج لم يتعلم جيدًا (ضعيف في التدريب والاختبار).
نختار ب: أعلى دقة اختبار ($83\%$) وفرق صغير.

\end{document}
